\section*{Capítulo \ref{ch:b3:organometallic-bonding} --- Química organometálica: enlace y ligandos}

\begin{solution}{exo:b3:organometallic-bonding:1}
$\ce{Fe(CO)5}$: $8 + 10 = 18$. $\ce{Mn2(CO)10}$: $7 + 10 + 1 = 18$ por Mn. $\ce{CpMo(CO)3H}$: $6 + 5 + 6 + 1 =
18$. $\ce{[PtCl3(C2H4)]-}$: $10 + 3 + 2 + 1 = 16$. $\ce{Cr($\eta^6$-C6H6)(CO)3}$: $6 + 6 + 6 = 18$. $\ce{Cp2ZrCl2}$:
$4 + 10 + 2 = 16$.
\end{solution}

\begin{solution}{exo:b3:organometallic-bonding:2}
Fe(0) $d^8$; Mn(0) $d^7$; Mo(II) $d^4$; Pt(II) $d^8$; Cr(0) $d^6$; Zr(IV) $d^0$.
\end{solution}

\begin{solution}{exo:b3:organometallic-bonding:3}
$\ce{[Mn(CO)6]+} > \ce{Cr(CO)6} > \ce{[V(CO)6]-}$: cuanto más rico en electrones es el metal, más
retrodona hacia $\pi^*$ y más débiles son los enlaces C--O.
\end{solution}

\begin{solution}{exo:b3:organometallic-bonding:4}
Etano ($\ce{Mn(CO)5}$ y $\ce{CpFe(CO)2}$ son \omterm{def:b3:organometallic-bonding:isolobal}{isolobulares} con el \ce{CH3}); de nuevo etano; tetraedrano,
$\ce{(CH)4}$, con tres CH sustituidos por $\ce{Co(CO)3}$.
\end{solution}

\begin{solution}{exo:b3:organometallic-bonding:5}
$fac$ ($C_{3v}$): dos (A$_1$ + E). $mer$ ($C_{2v}$): tres (2A$_1$ + B$_1$).
\end{solution}

\begin{solution}{exo:b3:organometallic-bonding:6}
La fosfina voluminosa congestiona el metal: la disociación alivia la tensión, de modo que es
más rápida, y se favorecen los intermedios de baja coordinación.
\end{solution}

\begin{solution}{exo:b3:organometallic-bonding:7}
El complejo de cromo es un \omterm{def:b3:organometallic-bonding:carbene}{carbeno de Fischer} (sustituyente OMe, Cr(0), ligandos CO): una
amina ataca el carbono carbénico electrófilo y sustituye al OMe. El alquilideno de
tántalo es un \omterm{def:b3:organometallic-bonding:carbene}{carbeno de Schrock}: su carbono nucleófilo ataca el carbonilo de una
cetona, lo que da un alqueno y una unidad Ta=O, como en una reacción de Wittig.
\end{solution}

\begin{solution}{exo:b3:organometallic-bonding:8}
Alternado: no hay solapamiento $\delta$, y los dos electrones $\delta$ son no enlazantes: orden de enlace 3. $\ce{[Re2Cl8]^{4-}}$:
$\sigma^2\pi^4\delta^2\delta^{\ast2}$, orden de enlace 3.
\end{solution}

\begin{solution}{exo:b3:organometallic-bonding:9}
Una distancia M$\cdots$H corta (difracción, mejor de neutrones); un número de onda de tensión C--H
bajo; una señal de RMN a campo alto con un acoplamiento C--H a un enlace reducido.
\end{solution}

\begin{solution}{exo:b3:organometallic-bonding:10}
16: los complejos plano-cuadrados $d^8$ ($\ce{[PtCl4]^{2-}}$, catalizador de Wilkinson), cuyo orbital del 18.º electrón
está demasiado alto; los \omterm{def:b3:organometallic-bonding:metallocene}{metalocenos} de los metales tempranos como el $\ce{Cp2ZrCl2}$, demasiado congestionados para admitir más
ligandos. 17: el $\ce{V(CO)6}$, que no puede dimerizarse por razones estéricas. 19: el cobaltoceno,
cuyo electrón adicional ocupa un orbital antienlazante (y se pierde con facilidad).
\end{solution}

\begin{solution}{exo:b3:organometallic-bonding:11}
El CO es un $\pi$-aceptor fuerte: baja el nivel $t_{2g}$ y hace $\Delta_{\mathrm o}$ grande, muy por encima de la
energía de apareamiento: $t_{2g}^6$, diamagnético. Las transiciones $d$--$d$, en $\Delta_{\mathrm o}$, caen en el ultravioleta,
y no se absorbe luz visible.
\end{solution}

\begin{solution}{exo:b3:organometallic-bonding:12}
La donación desde $\pi$ y la retrodonación hacia $\pi^*$ rebajan ambas el orden de enlace C--C. En la sal de
Zeise (Pt(II), retrodonación moderada) el enlace se alarga un poco; un metal de valencia baja,
rico en electrones, retrodona fuertemente y el complejo se acerca al límite del
\omterm{def:b3:organometallic-bonding:dcd}{metalaciclopropano}, con un enlace C--C próximo a un enlace sencillo.
\end{solution}

\begin{solution}{pb:b3:organometallic-bonding:1}
\textbf{1.} $\ce{Fe^{2+}}$ ($d^6$) $+ 2 \times 6 = 18$.
\textbf{2.} $8 + 2 \times 5 = 18$.
\textbf{3.} $9 + 10 = 19$.
\textbf{4.} $10 + 10 = 20$.
\textbf{5.} $18 - 1 = 17$.
\textbf{6.} Fe(II) $d^6$; Co(II) $d^7$; Ni(II) $d^8$; Fe(III) $d^5$.
\textbf{7.} $a_{1g}$ ($d_{z^2}$), $e_{1g}$ ($d_{xz}$, $d_{yz}$), $e_{2g}$ ($d_{xy}$, $d_{x^2-y^2}$).
\textbf{8.} El par $e_{1g}$ se solapa fuertemente con la combinación $e_{1g}$ llena de los anillos: orbitales
enlazantes situados sobre todo en los anillos, y antienlazantes $e_{1g}^\ast$, sobre todo en el metal, altos en
energía.
\textbf{9.} $e_{2g} \approx a_{1g} < e_{1g}^\ast$.
\textbf{10.} $e_{2g}^4a_{1g}^2$: ningún electrón desapareado.
\textbf{11.} Un electrón en $e_{1g}^\ast$: un electrón desapareado, $1.73\,\mu_B$.
\textbf{12.} Dos electrones en los dos orbitales $e_{1g}^\ast$, paralelos: dos desapareados, $\sqrt{8} =
2.83\,\mu_B$.
\textbf{13.} $e_{2g}^3a_{1g}^2$: un electrón desapareado. Un hueco $e_{2g}^3$ da un estado orbitalmente degenerado
(E), con \omterm{def:b3:complex-spectra-magnetism:moment}{contribución orbital}.
\textbf{14.} Su 19.º electrón está en un orbital antienlazante y se arranca con facilidad,
lo que da el catión cobaltocinio, de 18 electrones.
\textbf{15.} El electrón arrancado procede de un orbital casi no enlazante: la
estructura apenas cambia (\omterm{def:b3:complex-mechanisms:reorganisation}{energía de reorganización} pequeña, \cref{ch:b3:complex-mechanisms}).
\textbf{16.} El par es rápido y reversible en la mayoría de los disolventes, el compuesto es
soluble y estable, y su potencial depende poco del disolvente porque la
carga está repartida sobre un ion grande.
\textbf{17.} Con 20 electrones, dos de ellos antienlazantes, reacciona perdiéndolos:
adiciona ligandos con deslizamiento o pérdida de un anillo, o se oxida.
\textbf{18.} $\ce{CpFe(CO)2}$: $8 + 5 + 4 = 17$ electrones, a uno de 18, con un orbital frontera: isolobular
con el \ce{CH3}. El dímero $\ce{[CpFe(CO)2]2}$ es ``etano'', con un enlace Fe--Fe (en la práctica, dos ligandos CO
forman puentes sobre él).
\textbf{19.} $\ce{CpFe(CO)2-Mn(CO)5}$, con un enlace Fe--Mn.
\textbf{20.} La $\ce{PPh3}$ es mejor dador y peor $\pi$-aceptor que el CO: el metal retrodona
más hacia los CO restantes, cuyas tensiones bajan.
\textbf{21.} Dos (A$_1$ + E).
\textbf{22.} El deslizamiento del anillo a $\eta^3$ libera el equivalente de dos electrones de coordinación: un ligando
entrante puede unirse de forma asociativa sin que el metal supere los 18 electrones.
\textbf{23.} \textbf{$\mu_{\text{solo espín}}(\text{niqueloceno}) = \sqrt{2 \times 4} = 2.83\,\mu_B$.}
\end{solution}
